
\chapter{Método 4: Monte Carlo para Opciones Americanas}
\label{chapter:metodo_montecarlo}

La simulación de Monte Carlo es una técnica fundamental en finanzas computacionales que permite valorar derivados complejos mediante la generación de trayectorias aleatorias del precio del activo subyacente. Aunque es naturalmente adecuada para opciones europeas, su aplicación a opciones americanas presenta desafíos significativos debido a la posibilidad de ejercicio anticipado.

\section{Desafío de las Opciones Americanas}

\subsection{El Problema del Ejercicio Óptimo}

Las opciones americanas otorgan a su tenedor el derecho de ejercer en cualquier momento hasta el vencimiento $T$. Esta flexibilidad adicional implica que el valor de la opción en cualquier instante $t$ es:

\begin{equation}
V(S_t, t) = \max\left\{\text{Payoff inmediato}, \mathbb{E}^{\mathbb{Q}}[e^{-r\Delta t} V(S_{t+\Delta t}, t+\Delta t) | S_t]\right\}
\end{equation}

donde el primer término representa el valor de ejercicio inmediato y el segundo el valor de continuación (mantener la opción viva).

La decisión óptima de ejercicio requiere comparar en cada instante:
\begin{itemize}
    \item \textbf{Valor intrínseco}: $\max(S_t - K, 0)$ para calls o $\max(K - S_t, 0)$ para puts
    \item \textbf{Valor de continuación}: El valor esperado descontado de mantener la opción
\end{itemize}

\section{Simulación de Monte Carlo}

\subsection{Generación de Trayectorias}

Bajo el modelo de Black-Scholes, el precio del activo subyacente bajo la medida neutral al riesgo sigue:

\begin{equation}
dS_t = (r - \delta)S_t dt + \sigma S_t dW_t^{\mathbb{Q}}
\end{equation}

donde $r$ es la tasa libre de riesgo, $\delta$ es la tasa de dividendos, $\sigma$ es la volatilidad y $W_t^{\mathbb{Q}}$ es un proceso de Wiener bajo la medida neutral al riesgo.

La solución analítica de esta ecuación diferencial estocástica es:

\begin{equation}
S_{t+\Delta t} = S_t \exp\left[\left(r - \delta - \frac{\sigma^2}{2}\right)\Delta t + \sigma\sqrt{\Delta t}\, Z\right]
\end{equation}

donde $Z \sim N(0,1)$ es una variable aleatoria normal estándar.

\subsection{Algoritmo Básico para Opciones Europeas}

Para una opción europea, el algoritmo de Monte Carlo es directo:

\begin{enumerate}
    \item Generar $M$ trayectorias del precio del subyacente desde $t=0$ hasta $t=T$
    \item Para cada trayectoria $i$, calcular el payoff al vencimiento: $\text{Payoff}_i = \max(S_T^{(i)} - K, 0)$ para calls
    \item Estimar el valor de la opción como:
    \begin{equation}
    V(S_0, 0) \approx e^{-rT} \frac{1}{M} \sum_{i=1}^{M} \text{Payoff}_i
    \end{equation}
\end{enumerate}

El error estándar de esta estimación decrece como $O(1/\sqrt{M})$, independientemente de la dimensionalidad del problema.

\subsection{Monitoreo de Barreras}

Para opciones con barreras, es necesario monitorear si el precio cruza el nivel $H$ durante la vida de la opción. Esto requiere discretizar cada trayectoria en $N$ pasos temporales:

\begin{equation}
S_{t_j} = S_{t_{j-1}} \exp\left[\left(r - \delta - \frac{\sigma^2}{2}\right)\Delta t + \sigma\sqrt{\Delta t}\, Z_j\right]
\end{equation}

para $j = 1, 2, \ldots, N$, donde $\Delta t = T/N$ y $Z_j \sim N(0,1)$ son independientes.

En cada paso se verifica si $S_{t_j}$ ha cruzado la barrera $H$:
\begin{itemize}
    \item Para opciones \textbf{knock-out}: Si se cruza $H$, la trayectoria se marca como extinguida
    \item Para opciones \textbf{knock-in}: Si se cruza $H$, la trayectoria se marca como activada
\end{itemize}

\textbf{Sesgo de Discretización}: El monitoreo discreto puede no detectar cruces que ocurren entre puntos de observación. Este sesgo se reduce aumentando $N$, pero a costa de mayor tiempo computacional.

\subsection{El Problema Fundamental: ¿Cuándo Ejercer?}

Para opciones americanas, el desafío es determinar la estrategia de ejercicio óptima. En cada punto de decisión $t_j$, necesitamos saber si es óptimo ejercer o continuar.

El problema es que en Monte Carlo avanzamos hacia adelante en el tiempo, pero la decisión de ejercicio requiere conocer el valor de continuación, que depende del futuro. Esto crea una paradoja:

\begin{itemize}
    \item Para decidir si ejercer en $t_j$, necesitamos $\mathbb{E}[V(S_{t_{j+1}}, t_{j+1}) | S_{t_j}]$
    \item Pero $V(S_{t_{j+1}}, t_{j+1})$ depende de las decisiones de ejercicio futuras
    \item Que a su vez dependen de valores de continuación aún más futuros
\end{itemize}

Este problema circular requiere un enfoque sofisticado.

\section{Método de Longstaff-Schwartz (LSM)}

La simulación de Monte Carlo enfrenta una paradoja fundamental cuando se aplica a opciones americanas, derivada de la tensión entre la dirección temporal de la simulación y la naturaleza de la decisión de ejercicio óptimo. En Monte Carlo, generamos trayectorias aleatorias del precio del activo avanzando hacia adelante en el tiempo, desde $t=0$ hasta el vencimiento $T$. Sin embargo, para decidir si es óptimo ejercer la opción en un punto intermedio $t_j$, necesitamos comparar el valor de ejercicio inmediato (conocido) con el valor de continuación, que es el valor esperado de mantener la opción viva. Este valor de continuación depende de todas las decisiones de ejercicio futuras, que a su vez dependen de valores de continuación aún más futuros, creando una dependencia circular que no puede resolverse avanzando hacia adelante. El método de Longstaff-Schwartz resuelve esta paradoja mediante un truco estadístico elegante basado en regresión de mínimos cuadrados. La idea es generar primero todas las trayectorias completas hasta el vencimiento y luego trabajar hacia atrás en el tiempo (como en los árboles). En cada punto de decisión $t_j$, el método utiliza regresión para aproximar el valor de continuación como una función del precio actual del activo, típicamente usando polinomios de Laguerre como funciones base. Específicamente, para todas las trayectorias que están in-the-money en $t_j$, se realiza una regresión donde la variable independiente es el precio $S_{t_j}$ y la variable dependiente es el valor descontado que esas trayectorias alcanzaron en $t_{j+1}$ (conocido porque ya trabajamos hacia atrás desde el final). Los coeficientes de regresión resultantes proporcionan una estimación del valor de continuación condicional al precio actual, permitiendo compararlo con el valor de ejercicio inmediato y tomar la decisión óptima. Este enfoque convierte un problema de optimización dinámica intratable en una secuencia de problemas de regresión estándar, cada uno resoluble eficientemente mediante mínimos cuadrados ordinarios.

El método de \cite{longstaff2001valuing}, también conocido como Least Squares Monte Carlo (LSM), resuelve el problema del ejercicio óptimo mediante regresión para estimar el valor de continuación.

\subsection{Idea Central}

La intuición del método es:

\begin{enumerate}
    \item Generar todas las trayectorias completas desde $t=0$ hasta $t=T$
    \item Trabajar hacia atrás en el tiempo (backward induction)
    \item En cada punto de decisión $t_j$, usar regresión para estimar el valor de continuación como función del precio actual
    \item Comparar el valor de continuación estimado con el valor de ejercicio inmediato
    \item Ejercer si el valor inmediato excede el valor de continuación
\end{enumerate}

\subsection{Estimación del Valor de Continuación}

En el tiempo $t_j$, para todas las trayectorias que están ``in-the-money'' (donde el ejercicio inmediato tiene valor positivo) se estima el valor de continuación mediante regresión de mínimos cuadrados \textbf{sobre los flujos de caja reales obtenidos ex-post} (y no sobre el valor esperado de la opción, lo cual generaría un sesgo alcista):

\begin{equation}
\mathbb{E}[\sum_{i=j+1}^{N}  e^{-r(t_i - t_j)} CF_i | S_{t_j}] \approx \sum_{k=0}^{K} a_k L_k(S_{t_j})
\end{equation}

donde $CF_i$ representa el flujo de caja real que ocurre en el tiempo $t_i$ a lo largo de la trayectoria condicionado a que la opción no se ejerció antes, $L_k$ son funciones base (típicamente polinomios de Laguerre) y $a_k$ son coeficientes determinados por regresión.

\textbf{¿Por qué Polinomios de Laguerre?}

Los polinomios de Laguerre ponderados $L_k(x) e^{-x/2}$ forman una base ortogonal en $[0, \infty)$, que es el dominio natural de los precios de activos. Por simplicidad se utilizan polinomios de Laguerre simples (no ponderados) ya que la elección exacta de las funciones base (ya sean ponderadas, no ponderadas, de Hermite, etc.) no afecta significativamente la convergencia del algoritmo siempre que se usen los términos suficientes. Los primeros polinomios son:

\begin{align}
L_0(x) &= 1 \\
L_1(x) &= 1 - x \\
L_2(x) &= 1 - 2x + \frac{x^2}{2} \\
L_3(x) &= 1 - 3x + \frac{3x^2}{2} - \frac{x^3}{6}
\end{align}

En la práctica, se usan típicamente los primeros 3-5 polinomios.

\subsection{Algoritmo de Longstaff-Schwartz}

El algoritmo LSM procede en cuatro pasos principales:

\textbf{Paso 1: Generación de Trayectorias}

Se generan $M$ trayectorias completas del precio del subyacente desde $t=0$ hasta $t=T$, discretizadas en $N$ pasos temporales.

\textbf{Paso 2: Inicialización al Vencimiento}

Para cada trayectoria, el valor al vencimiento es simplemente el payoff de la opción.

\textbf{Paso 3: Inducción Hacia Atrás}

Este es el núcleo del algoritmo. Para cada tiempo $t_j$ desde $t_{N-1}$ hasta $t_1$:

\begin{enumerate}
    \item Se identifican las trayectorias in-the-money (donde el ejercicio inmediato tiene valor positivo)
    
    \item Para estas trayectorias, se realiza una regresión de mínimos cuadrados usando funciones base (típicamente polinomios de Laguerre) para estimar el valor de continuación:
    \begin{equation}
    \mathbb{E}[e^{-r\Delta t} V(S_{t_{j+1}}, t_{j+1}) | S_{t_j}] \approx \sum_{k=0}^{K} a_k L_k(S_{t_j})
    \end{equation}
    
    \item Se compara el valor de ejercicio inmediato con el valor de continuación estimado
    
    \item Se toma la decisión de ejercicio óptima: ejercer si el valor inmediato excede el valor de continuación
\end{enumerate}

\textbf{Paso 4: Valoración}

El valor de la opción es el promedio descontado de los payoffs en los tiempos de ejercicio óptimos determinados en el paso anterior.

\subsection{Extensión a Opciones Barrera Americanas}

Para opciones americanas con barreras, el algoritmo LSM se modifica para incorporar el monitoreo de la barrera:

\begin{itemize}
    \item \textbf{Knock-out}: Si una trayectoria cruza la barrera en $t_j$, se establece su valor en el rebate $R$ y no se consideran decisiones de ejercicio futuras para esa trayectoria.
    
    \item \textbf{Knock-in}: La trayectoria solo puede ejercerse después de que se haya cruzado la barrera. Antes del cruce, solo se considera el valor de continuación.
\end{itemize}

\subsection{Ventajas y Limitaciones del Método LSM}

\textbf{Ventajas:}
\begin{itemize}
    \item \textbf{Flexibilidad}: Aplicable a opciones con payoffs complejos, múltiples subyacentes y barreras
    \item \textbf{Escalabilidad}: No sufre de la maldición de la dimensionalidad como árboles y diferencias finitas
    \item \textbf{Implementación}: Relativamente simple de programar
    \item \textbf{Griegas}: Pueden estimarse mediante diferencias finitas o métodos de sensibilidad
\end{itemize}

\textbf{Limitaciones:}
\begin{itemize}
    \item \textbf{Convergencia lenta}: Error $O(1/\sqrt{M})$ requiere muchas simulaciones para alta precisión
    \item \textbf{Sesgo bajo}: El método tiende a subestimar ligeramente el valor verdadero (sesgo bajo)
    \item \textbf{Elección de funciones base}: El rendimiento depende de la elección apropiada de funciones base
    \item \textbf{Costo computacional}: Para alta precisión, puede ser más lento que métodos de árbol para problemas de baja dimensionalidad
\end{itemize}

\subsection{Comparación con Otros Métodos}

El método LSM es particularmente valioso cuando:
\begin{itemize}
    \item El problema es de alta dimensionalidad (opciones sobre múltiples activos)
    \item El payoff es complejo o path-dependent
    \item Se requiere flexibilidad para incorporar características exóticas
\end{itemize}

Para opciones americanas simples de baja dimensionalidad, los árboles adaptativos (como el AMM) pueden ser más eficientes. Sin embargo, para opciones barrera americanas con múltiples características exóticas, el LSM es frecuentemente el método de elección en la práctica.
